\documentclass[11pt]{article}
\usepackage[a4paper,margin=28mm]{geometry}
\usepackage[T1]{fontenc}
\usepackage{lmodern}
\usepackage{amsmath,amssymb,amsthm}
\usepackage{microtype}
\usepackage[colorlinks=true,linkcolor=blue,citecolor=blue,urlcolor=blue]{hyperref}
\newtheorem{theorem}{Theorem}
\newtheorem{corollary}[theorem]{Corollary}
\newcommand{\meet}{\wedge}
\newcommand{\dotop}{\mathbin{\cdot}}
\title{Semibrace existence on meet-semilattices:\\
an attributed classification}
\author{Alec Kriebel\\
\small Independent researcher\\
\small ORCID: \href{https://orcid.org/0009-0001-9320-500X}{0009-0001-9320-500X}}
\date{17 October 2026}

\begin{document}
\maketitle

\begin{abstract}
We give a self-contained verification of the two existence questions in
Dietzel's Problem~3, under the axioms printed in Oberwolfach Report~9/2023.
Every nonempty meet-semilattice admits the printed semibrace structure.
A simultaneous semibrace and \(L\)-algebra structure exists exactly when the
semilattice has a greatest element; in every such structure the two orders
coincide. The operation used for sufficiency is exactly the operation of
Castiglioni and San Mart\'in, arXiv:1611.09589v1 (2016), Example~3.5(4).
This note records an attributed verification and deduction of the full
historical existence answer, without a claim of a new solution, operation,
proof mechanism, or priority. We distinguish a separate convention requiring
the distinguished element to be a meet identity.
\end{abstract}

\section{The printed questions and their conventions}

Let \((X,\meet)\) be a meet-semilattice: \(\meet\) is associative,
commutative, and idempotent. Write \(x\leq y\) when \(x\meet y=x\).
We allow the empty carrier solely to state the boundary case explicitly.
Dietzel's Definition~2 and Problem~3 \cite[p.~545]{Dietzel} ask for a
point \(e\in X\) and a binary operation \(\dotop\) satisfying, for all
\(x,y,z\in X\),
\begin{align}
x\dotop(y\meet z)&=(x\dotop y)\meet(x\dotop z), \tag{S1}\label{S1}\\
e\dotop x&=x, \tag{S2}\label{S2}\\
(x\meet y)\dotop z&=(x\dotop y)\dotop(x\dotop z), \tag{S3}\label{S3}\\
x\dotop e&=e. \tag{S4}\label{S4}
\end{align}
The semilattice footnote on the preceding page specifies the three meet
laws above; no meet identity is included in Definition~2.
The second question requires the same structure also to be an
\(L\)-algebra. In addition to \eqref{S2} and \eqref{S4}, Definition~1
\cite[p.~544]{Dietzel} requires
\begin{align}
x\dotop x&=e, \tag{L1}\label{L1}\\
(x\dotop y)\dotop(x\dotop z)&=(y\dotop x)\dotop(y\dotop z), \tag{L2}\label{L2}\\
x\dotop y=e=y\dotop x&\quad\Longleftrightarrow\quad x=y. \tag{L3}\label{L3}
\end{align}
Equation~\eqref{L2} is the cycloid identity.

There is a relevant difference of source conventions. Definition~3.1 in
Dietzel, Mench\'on, and Vendramin \cite{Enumeration} first requires
\((X,\meet,e)\) to be a commutative monoid and then imposes the same four
equations. On a meet-semilattice this adds \(x\meet e=x\), equivalently that
\(e\) is greatest. We keep these printed definitions separate. This
comparison does not assert which additional axioms an author may have
intended.
The enumeration paper attributes the semibrace definition and closure
framework to Rump. Its version~v1, dated 10 June 2022, already prints the
monoid convention in Definition~3.1; this convention predates the workshop.

\begin{theorem}\label{main}
For the workshop definition \eqref{S1}--\eqref{S4}:
\begin{enumerate}
\item A meet-semilattice admits a semibrace structure if and only if it is
nonempty. Every prescribed point \(e\in X\) is allowed.
\item A meet-semilattice admits a simultaneous semibrace and \(L\)-algebra
structure if and only if it has a greatest element. For a prescribed point
\(e\), this is possible exactly when \(e\) is that greatest element. In
every simultaneous structure,
\[
x\meet y=x\quad\Longleftrightarrow\quad x\dotop y=e.
\]
\end{enumerate}
\end{theorem}

The sufficient operation for the second assertion is the historical formula
\begin{equation}\label{construction}
x\dotop y=
\begin{cases}
e,&x\leq y,\\
x\meet y,&x\not\leq y.
\end{cases}
\end{equation}
It appears, with \(1\) in place of \(e\), in Castiglioni and San Mart\'in
\cite[p.~7, Example~3.5(4)]{CastiglioniPreprint}. Their term
\emph{bounded semilattice} means one with a greatest element
\cite[p.~2]{CastiglioniPreprint}. The following verification uses that
inspected 2016 text for the attribution of the operation. The later,
retitled publication is listed separately in \cite{CastiglioniPublished};
its publication metadata are distinct from the versioned text inspected
here.

\section{Proof of the existence classification}

For the first assertion, on a nonempty \(X\) with any point \(e\), set
\(x\dotop y=y\). Both sides of \eqref{S1} are \(y\meet z\), and both sides of
\eqref{S3} are \(z\); \eqref{S2} and \eqref{S4} are immediate. An empty set
has no point \(e\). This projection is an \(L\)-algebra only on a singleton,
since its diagonal is \(x\dotop x=x\).

\medskip
\noindent\emph{Necessity for the simultaneous structure.}
Suppose all the displayed axioms hold. Fix \(x\) and put \(a=x\meet e\).
Twice applying \eqref{S3} gives
\[
a\dotop x=(x\dotop e)\dotop(x\dotop x)=e\dotop e=e,
\qquad
a\dotop a=(x\dotop e)\dotop(x\dotop a)=x\dotop a.
\]
By \eqref{L1}, \(a\dotop a=e\), so \(x\dotop a=e=a\dotop x\).
Separation \eqref{L3} yields \(a=x\). Thus \(x\meet e=x\) for every \(x\);
\(e\) is greatest.

The order equality follows by equally local substitutions. If \(x\meet y=x\),
then \eqref{S3} with \(z=y\), followed by \eqref{L1}, gives
\[
x\dotop y=(x\dotop y)\dotop(x\dotop y)=e.
\]
Conversely, suppose \(x\dotop y=e\) and put \(a=x\meet y\). Now
\[
a\dotop x=(x\dotop y)\dotop(x\dotop x)=e,
\qquad
a\dotop a=(x\dotop y)\dotop(x\dotop a)=x\dotop a.
\]
The diagonal and separation again imply \(a=x\). No injectivity of an
individual left translation is used.

\medskip
\noindent\emph{Sufficiency.}
Let \(e\) be greatest and define \(\dotop\) by \eqref{construction}.
Since \(x\leq x,e\), the diagonal and \eqref{S4} hold.
If \(y=e\), then \(e\dotop y=e=y\); otherwise \(e\not\leq y\) and
\(e\dotop y=e\meet y=y\). This proves \eqref{S2}.
In the second branch of \eqref{construction}, \(x\meet y=e\) would force
\(x=y=e\), contrary to \(x\not\leq y\). Hence
\[
x\dotop y=e\quad\Longleftrightarrow\quad x\leq y.
\]
Antisymmetry of the meet order proves \eqref{L3}.

To prove meet preservation \eqref{S1}, use
\(x\leq y\meet z\) if and only if \(x\leq y\) and \(x\leq z\).
The four exhaustive cases give the following common values for its two
sides:
\[
\begin{array}{c|c|c}
x\leq y&x\leq z&\text{common value}\\ \hline
\text{yes}&\text{yes}&e\\
\text{yes}&\text{no}&x\meet z\\
\text{no}&\text{yes}&x\meet y\\
\text{no}&\text{no}&x\meet y\meet z
\end{array}
\]
For example, in the second row the left side is
\(x\meet y\meet z=x\meet z\), while the right side is
\(e\meet(x\meet z)=x\meet z\). The third row is its symmetric counterpart;
the first and fourth follow directly from the formula and the meet laws.

For \eqref{S3}, if \(x\leq y\), its right side is
\(e\dotop(x\dotop z)=x\dotop z=(x\meet y)\dotop z\).
Otherwise put \(a=x\meet y\); then \(x\dotop y=a\).
If \(x\leq z\), both \(a\dotop z\) and \(a\dotop(x\dotop z)=a\dotop e\)
are \(e\). If \(x\not\leq z\), put \(b=x\meet z\), so \(x\dotop z=b\).
Because \(a\leq x\), we have
\[
a\leq b\quad\Longleftrightarrow\quad a\leq z.
\]
When these conditions hold, \(a\dotop b=a\dotop z=e\).
When they fail, \(a\dotop b=a\meet b=a\meet z=a\dotop z\).
Thus \eqref{S3} holds in every branch. Finally, commutativity of meet gives
\[
(x\dotop y)\dotop(x\dotop z)
=(x\meet y)\dotop z
=(y\meet x)\dotop z
=(y\dotop x)\dotop(y\dotop z),
\]
which is \eqref{L2}. This proves Theorem~\ref{main}.

\begin{corollary}
If the semibrace definition additionally requires \(e\) to be a meet
identity, both existence questions have exactly the meet-semilattices with
a greatest element as their answer.
\end{corollary}
\begin{proof}
A meet identity is a greatest element. When it exists, projection proves
the first existence assertion and \eqref{construction} proves the second.
\end{proof}

\section{Boundary cases, scope, and status}

The singleton satisfies both questions, with its unique point and operation.
The empty carrier satisfies neither pointed existence question.
No finiteness, least element, joins, completeness, or distributivity is
assumed. For example, \(\mathbb Z_{\leq0}\) with meet given by minimum has
greatest element \(0\) and no least element, and \eqref{construction}
applies. It also applies to the nondistributive five-element diamond
\(M_3\). The proof uses only binary meets, their elementary order laws, and
the greatest element.

The formula with \(y\), rather than \(x\meet y\), in its second branch is
a different operation, displayed in Dietzel's Problem~10
\cite[p.~546]{Dietzel}. It must not be substituted for
\eqref{construction}: in the Boolean square with distinct atoms \(a,b\),
it gives \(a\dotop(a\meet b)=0\), but
\((a\dotop a)\meet(a\dotop b)=e\meet b=b\).
This concerns meet preservation and does not dispute the separate
\(L\)-algebra assertion made there.

Theorem~\ref{main} classifies the semilattices admitting at least one
multiplication; it does not classify all compatible multiplications.
It makes no assertion of self-similarity or of a structure involving a
group operation. It answers the historical existence questions under the
explicit printed axioms. No assertion is made about current worldwide
openness, independent discovery, a first combined statement, or historical
priority for the deduction or its proof.
The full 2008 papers underlying Rump's framework were unavailable to the
current audit; bounded searches do not exclude earlier explicit or implicit
resolutions.

\paragraph{Preparation and verification disclosure.}
AI tools were used extensively in solving, reconstructing, checking,
prose and code preparation, and literature work leading to this note.
The manuscript is unrefereed. No external human review or formal
certification is claimed. The mathematical argument is the explicit
proof above, rather than a finite computation or a review label.

\begin{thebibliography}{9}
\bibitem{Dietzel}
C.~Dietzel, \emph{Some problems on \(L\)-algebras}, in
\emph{Mini-Workshop: Skew Braces and the Yang--Baxter Equation},
Oberwolfach Report 9/2023, pp.~544--547.
\href{https://ems.press/content/serial-article-files/47003}{Official EMS report};
\href{https://doi.org/10.4171/OWR/2023/9}{doi:10.4171/OWR/2023/9}.

\bibitem{Enumeration}
C.~Dietzel, P.~Mench\'on, and L.~Vendramin,
\emph{On the enumeration of finite \(L\)-algebras},
\href{https://arxiv.org/abs/2206.04955v4}{arXiv:2206.04955v4},
2 July 2023, Definition~3.1, p.~5.
The same convention occurs in
\href{https://arxiv.org/abs/2206.04955v1}{version~v1}, 10 June 2022,
Definition~3.1, p.~5.

\bibitem{CastiglioniPreprint}
J.~L.~Castiglioni and H.~J.~San Mart\'in,
\emph{On a symmetrization of hemiimplicative semilattices},
\href{https://arxiv.org/abs/1611.09589v1}{arXiv:1611.09589v1},
29 November 2016, pp.~2 and~7, Example~3.5(4).

\bibitem{CastiglioniPublished}
J.~L.~Castiglioni and H.~J.~San Mart\'in,
\emph{\(\ell\)-Hemi-Implicative Semilattices},
Studia Logica \textbf{106} (2018), 675--690.
\href{https://doi.org/10.1007/s11225-017-9759-3}{doi:10.1007/s11225-017-9759-3}.
Published online 3 October 2017. Retitled publication of the preceding
preprint; only the versioned preprint body is used for the formula
attribution in this note.
\end{thebibliography}
\end{document}
